\chapter{Informação, Cultura e Poder: Quem Decide o que Parece Verdade}
\label{ch:poder}

Considera uma situação fictícia. Recebes uma gravação curta, acompanhada de uma mensagem urgente: a escola fecha amanhã. A voz parece conhecida. A gravação pode ser fabricada, antiga, cortada ou inteiramente verdadeira. O ficheiro, por si só, ainda não resolveu a pergunta que te interessa.

A investigação começa por separar duas coisas: o objeto que chegou e a alegação que o acompanha. Uma fotografia verdadeira pode ter uma legenda falsa. Um artigo real pode ser citado para sustentar uma frase que nunca demonstrou. A IA generativa amplia os recursos de produção; a confiança continua a depender de origem, contexto e confirmação.

\section{Deepfakes e a diferença entre aparência e acontecimento}
Um deepfake pode alterar rosto, voz ou ação para produzir uma representação enganosa. Nem toda a edição é enganosa: cinema, sátira identificada e acessibilidade usam transformação de media há muito tempo. Finalidade, apresentação, consentimento e consequências fazem diferença.

Para verificar uma peça, procura a publicação de origem, contexto, data e confirmação independente. Uma chamada urgente com voz familiar não deve alterar sozinha um procedimento de confirmação. A resposta organizacional mais robusta usa um canal já conhecido e uma regra prévia, em vez de confiar apenas na capacidade humana de detetar artefactos.

Detetores podem apoiar triagem, mas o seu alcance depende dos dados e transformações avaliados. Uma compressão, captura de ecrã ou ferramenta nova pode alterar desempenho. Um score isolado não deve servir como prova definitiva de autenticidade ou fraude.

\section{A pergunta muda o que conta como prova}
«Este áudio foi gerado?» é uma pergunta sobre produção. «A escola fecha amanhã?» é uma pergunta sobre um acontecimento. Mesmo que resolvêssemos a primeira, ainda poderíamos precisar de confirmar a segunda. Uma voz verdadeira pode repetir uma informação errada; uma voz sintética pode ler um aviso verdadeiro e devidamente identificado.

No cenário fictício, o passo útil é procurar o aviso no canal oficial já conhecido ou confirmar por um contacto obtido de forma independente da mensagem. Reencaminhar a gravação para vinte grupos e esperar que alguém saiba pode aumentar precisamente a circulação que estamos a tentar avaliar. O procedimento precisa de existir antes da urgência, quando ainda temos tempo para o combinar.

Esta disciplina também evita atribuir valor excessivo a pormenores visuais. Uma mão estranha pode sugerir geração; não prova a falsidade de toda a alegação. Uma mão perfeita não prova autenticidade. O que procuramos é uma ligação verificável entre a afirmação e um acontecimento ou documento, não uma competição para encontrar o píxel mais suspeito.

Há conclusões diferentes a registar. «Confirmado» requer suporte adequado. «Contradito» pede evidência contrária. «Não verificado» conserva uma lacuna. A terceira categoria pode ser frustrante numa conversa rápida, mas evita que a ausência de informação se transforme numa certeza que ninguém encontrou.

\section{Proveniência criptográfica tem um alcance específico}
As Content Credentials da C2PA descrevem proveniência e ligações verificáveis entre conteúdo e declarações~\citep{c2pa22}. Podem ajudar a verificar integridade e história registada. Não certificam a verdade de todas as interpretações sobre uma cena.

Uma imagem assinada pode mostrar uma encenação; uma imagem sem metadados pode ser verdadeira. A cadeia depende também de quem emite declarações e do que ficou fora do registo. Confiança digital pede uma combinação de mecanismos, fontes e instituições, não um botão que elimina o julgamento.

Podemos imaginar um recibo da história de um ficheiro: quem declarou uma operação, que ligação existe entre versões e se a integridade dessa ligação pode ser verificada. A analogia termina antes de «o recibo prova a verdade». O mecanismo regista declarações dentro de uma cadeia de confiança; não esteve em todos os locais que uma imagem pretende representar.

Uma verificação útil conserva, por isso, perguntas separadas. A cadeia é válida segundo o mecanismo usado? Quem fez as declarações relevantes? O que dizem exatamente? Que interpretação do acontecimento precisa de confirmação fora dessa cadeia? Proveniência acrescenta informação ao processo, sem absorver a função inteira de verificar.

\section{Recomendação e distribuição são parte do problema}
Produzir conteúdo e decidir quem o vê são funções diferentes. Um sistema de recomendação pode afetar alcance, incentivos e concentração de atenção mesmo sem gerar uma frase. Uma análise do impacto da IA deve incluir estes sistemas, em vez de reduzir tudo a chatbots.

Num cenário conceptual, conteúdo barato aumenta o volume disponível. A atenção humana e a capacidade de verificação continuam limitadas. Pode haver mais variedade, mais repetição ou ambas em grupos diferentes. Para afirmar um efeito sobre opinião pública, precisamos de medir exposição, mudança e alternativas; volume de publicações não demonstra persuasão por si só.

No jornalismo, um resumo automático pode ajudar a organizar documentos. O processo precisa de conservar fontes, confirmar alegações e distinguir material gerado de reportagem realizada. Entrevistas inventadas e citações fabricadas não se tornam investigação por aparecerem num estilo jornalístico.

Um resumo também pode deslocar o sentido sem inventar uma frase inteira. Se retira a condição «nas tarefas avaliadas» ou troca «associado» por «causado», a notícia fica mais forte do que o estudo. Uma revisão factual deve comparar a afirmação completa, incluindo qualificações e denominadores, com a fonte. Encontrar o título do artigo não basta.

\section{Persuasão não é uma medida de verdade}
Salvi e colaboradores estudaram debates estruturados com novecentos participantes e condições com humanos ou GPT-4, com e sem acesso a informação pessoal. A condição IA com personalização superou a base humano--humano na medida estudada. A comparação direta entre as duas condições IA não atingiu significância estatística, com $p=0.07$~\citep{salvi2025}. O resultado não permite atribuir toda a diferença à personalização.

O desenho mede mudança de posição em debates controlados, não votos numa eleição ou comportamento duradouro em redes sociais. Também não mede a verdade de uma opinião. É um caso útil para separar capacidade de persuadir, efeito em condições reais e justificação de uma alegação. Nenhuma dessas perguntas deve ser substituída pelas outras.

No desenho de um serviço, a diferença pode orientar limites de finalidade. Personalizar uma explicação para torná-la acessível e explorar informação íntima para pressionar uma decisão são usos distintos. Avaliá-los exige conhecer o objetivo, os dados usados, a possibilidade de recusar e a autonomia de quem recebe a mensagem. A posição editorial desta obra é tornar essa influência discutível e visível, em vez de a esconder sob a palavra «ajuda».

\section{Cultura, autoria e trabalho criativo}
Há perguntas distintas: de onde vieram os materiais de treino; que direitos permitem o seu uso; quem controla uma criação; que condições se aplicam à distribuição; e como o rendimento é repartido. As respostas jurídicas dependem da jurisdição e do caso. A autorização de uma plataforma não resolve automaticamente direitos de terceiros.

Para um projeto, conserva as referências utilizadas, licenças, intervenções humanas e termos relevantes. Se usares materiais licenciados com restrições, essas condições não desaparecem porque houve transformação por IA. O capítulo~\ref{ch:imagem-design} aplica este cuidado a um cartaz; o mesmo princípio orienta recursos digitais da obra.

Também existe uma decisão editorial sobre valor cultural que não é apenas uma questão de direitos. Podemos querer diversidade de vozes, remuneração e capacidade de participação. São objetivos a discutir explicitamente. Uma média de avaliações de qualidade não mede todos esses objetivos.

Um arquivo de criação ajuda a separar as contribuições. Guarda o material de partida autorizado, o pedido, as versões escolhidas e as alterações humanas relevantes. Pode não resolver uma questão jurídica, mas permite descrever o processo sem inventar uma origem. Para um leitor ou um cliente, essa descrição pode ser mais informativa do que o rótulo binário «feito por IA».

Há ainda escolhas culturais que uma métrica de produtividade não decide. Uma escola pode querer que os alunos experimentem escrever; uma comunidade pode valorizar registos locais; um artista pode procurar um processo lento por razões expressivas. Tornar a produção mais rápida é uma capacidade. Escolher quando a usar envolve finalidades que merecem ser ditas.

\section{Justiça e serviços públicos: contestar uma decisão}
Num serviço público conceptual, a IA pode ajudar a localizar regras ou preparar uma minuta. Selecionar pessoas para fiscalização, priorizar pedidos ou recomendar uma decisão afeta direitos e oportunidades. Finalidade, competência institucional e obrigações aplicáveis precisam de análise própria, como explica o capítulo~\ref{ch:governacao}.

Uma supervisão útil permite perceber dados, corrigir erro e alterar o resultado. Se o cidadão só recebe «o sistema decidiu», a automatização tornou a decisão menos contestável. Uma explicação gerada posteriormente não substitui o registo do processo efetivo.

Avalia também falsos positivos, grupos afetados, qualidade dos dados históricos e ciclos de retroalimentação. Se fiscalizamos mais um grupo, encontraremos mais registos nesse grupo; treinar sobre esses registos sem considerar a política de recolha pode reforçar a mesma seleção. O exemplo descreve um mecanismo possível, cuja presença num caso concreto deve ser demonstrada.

\section{SyRI, o que acontece quando o Estado assinala um risco}
Em fevereiro de 2020, um tribunal neerlandês teve de avaliar uma ferramenta concebida para detetar fraude. O \emph{Systeem Risico Indicatie}, ou SyRI, cruzava dados já detidos por entidades públicas. Um modelo com indicadores de risco ajudava a selecionar casos. Após análise adicional, podia ser emitido um relatório de risco para eventual investigação~\citep[§§ 4.27--4.29 e 6.62]{syri2020}.

O relatório podia ficar num registo durante dois anos sem que a pessoa afetada soubesse. O tribunal considerou significativo esse efeito na vida privada. Uma sanção teria de assentar também noutros elementos. A análise judicial incidiu no artigo 8.º, n.º 2, da Convenção Europeia dos Direitos Humanos. A classificação jurídica do processo ao abrigo do artigo 22.º do RGPD ficou por decidir nesse processo~\citep[§§ 6.54--6.60]{syri2020}.

O Tribunal Distrital de Haia examinou o equilíbrio entre combater fraude e respeitar a vida privada. Entre os problemas, identificou a falta de informação verificável sobre os indicadores e o modelo de risco, além de garantias insuficientes quanto ao volume e à necessidade dos dados. Em 5 de fevereiro de 2020, declarou sem efeito vinculativo, relativamente às organizações indicadas e às pessoas que representavam, o artigo 65.º da lei SUWI e o capítulo 5a do decreto SUWI~\citep[§§ 6.49, 6.80, 6.106 e 7.2]{syri2020}.

Para desenhar um serviço semelhante, seguiríamos o percurso de uma sinalização. Quem autorizou o cruzamento? Que dado sustentou a análise? Quem pode verificar a versão da regra e corrigir um erro? A evidência para responder a cada pergunta é diferente. Uma taxa de acerto vem de um conjunto avaliado. A possibilidade de contestar exige um processo que a pessoa consiga usar. O caso SyRI permite discutir este segundo problema mesmo sem estimar a precisão do modelo.

\section{Dois critérios de equidade podem pedir coisas diferentes}
Uma discussão sobre enviesamento precisa de especificar o erro e a comparação. Queremos a mesma taxa de falsos positivos? A mesma capacidade de detetar casos? A mesma proporção de acertos entre as pessoas sinalizadas? Ou acesso equivalente a uma revisão? Não são maneiras diferentes de escrever uma única medida.

O trabalho de Chouldechova analisa conflitos entre critérios estatísticos em instrumentos de previsão de reincidência, incluindo o papel de diferenças de prevalência~\citep{chouldechova2017}. É um enquadramento para tornar as escolhas explícitas, não uma decisão automática sobre justiça. Os pressupostos matemáticos e o contexto da decisão têm de acompanhar qualquer conclusão.

Podemos ver uma diferença sem usar dados de pessoas. Os dois cenários sintéticos do capítulo~\ref{ch:ciencia} conservam sensibilidade de 90\% e especificidade de 95\%. Com prevalências impostas de 1\% e 10\%, os valores preditivos positivos são cerca de 15.4\% e 66.7\%. São grupos matemáticos construídos, sem correspondência com populações sociais. Igualar duas taxas não igualou todas as outras grandezas.

Isto não prova que uma disparidade real é aceitável ou inevitável. Mostra por que precisamos de definir o critério, examinar a origem dos dados e discutir consequências. Num caso real, uma diferença pode resultar da recolha, da etiqueta, da seleção dos casos ou de outros mecanismos. Uma tabela agregada não identifica sozinha a causa, tal como remover um atributo explícito não elimina informação correlacionada noutros campos.

Uma auditoria útil acompanha o percurso inteiro: quem consegue apresentar o pedido, que dados entram, que saída recebe, que ação se segue e como pode corrigir um erro. Melhorar uma métrica no fim não compensa necessariamente uma barreira de acesso no início. A decisão de implementação deve explicar que valores protege e que trade-offs aceita, dentro das obrigações aplicáveis.

\section{A pessoa no circuito precisa de poder mudar o resultado}
Num serviço público imaginado, o sistema prepara uma recomendação e uma pessoa carrega em «validar». Isso demonstra a existência de um clique humano. Para saber se há supervisão efetiva, precisamos de examinar o tempo disponível, a informação acessível e a autoridade para contrariar a proposta.

Uma revisão pode falhar se a pessoa só vir a conclusão e uma explicação gerada posteriormente. O registo útil deve permitir reconstruir os dados e regras relevantes, dentro das permissões e da minimização adequada. Se a falha veio de um documento antigo, interessa identificar essa versão; se veio de um dado incorreto, interessa corrigi-lo e rever o efeito sobre a decisão.

O caminho de contestação faz parte do serviço. Quem pede revisão precisa de saber como fazê-lo e de conseguir apresentar informação que o sistema não considerou. Publicar uma descrição técnica incompreensível pode cumprir uma intenção de transparência sem tornar a decisão mais contestável. A proposta de desenho é avaliar o que a pessoa consegue efetivamente perceber e fazer.

\section{Vigilância e assimetria de informação}
Combinar dados permite inferências que um ficheiro isolado não revelava. Reconhecer padrões de deslocação, comportamento ou relações pode ter usos legítimos e abusivos. A disponibilidade técnica não elimina necessidade, proporcionalidade, base jurídica e controlo institucional.

Uma pergunta de projeto é quem ganha capacidade de observar e quem ganha capacidade de recusar ou contestar. A transparência precisa de ser útil para a pessoa afetada: saber que há IA não chega se ninguém explica o efeito da decisão ou o caminho de revisão.

No nosso cenário da gravação, a investigação termina quando a alegação encontra confirmação independente, é contradita ou permanece sem suporte suficiente para agir. Nos sistemas institucionais, a mesma disciplina tem de sobreviver a escala: conservar a relação entre afirmação, fonte e decisão, e permitir que alguém a examine. O poder de gerar informação aumenta a importância de poder verificar e contestar o seu uso.

\begin{promptlabbox}{Verificar a peça antes de a partilhar}
\textbf{Objetivo e contexto:} rever uma notícia ou imagem pública sem a redistribuir. Usa uma ficha com alegação, origem, data, confirmação e dúvida.

\textbf{Procedimento:} localiza a fonte original; confirma se os media correspondem à data e ao acontecimento; procura uma fonte independente; regista o que credenciais ou detetores realmente permitem concluir.

\textbf{Resultado esperado e validação:} uma decisão justificada entre confirmado, falso ou não verificado. Ausência de metadados não é classificada automaticamente como falsidade.

\textbf{Falhas possíveis:} partilhar para «perguntar se é verdade», aceitar uma voz familiar como autenticação ou confundir uma assinatura com prova de um acontecimento.
\end{promptlabbox}

Até aqui olhámos sobretudo para a informação que recebemos. No próximo capítulo, a conversa aproxima-se: o sistema lembra preferências, responde com paciência e parece acompanhar-nos. A pergunta deixa de ser apenas se o conteúdo é verdadeiro. Passa também por compreender o que estamos a aprender e a sentir nessa interação.
